
est aussi repr\'esentable par une immersion ferm\'ee.}
\item \emph{@@M0@@ est de type fini sur @@M1@@, c'est-\`a-dire, la famille des classes des faisceaux inversibles @@M2@@-\'equivalents \`a @@M3@@ sur les fibres de @@M4@@ est limit\'ee.}
\end{enumerate}

En effet, dans (i) (resp. (ii)), on doit traiter le cas d'un morphisme @@M5@@, o\`u @@M6@@ est un @@M7@@-sch\'ema quelconque, et montrer que l'image inverse de @@M8@@ est repr\'esentable par un sous-sch\'ema ouvert (resp. et ferm\'e) de @@M9@@. Il suffit de prendre @@M10@@ sur un ouvert affine @@M11@@ de @@M12@@. Alors @@M13@@ est la limite du syst\`eme projectif filtrant des @@M14@@-sch\'emas affines @@M15@@, o\`u @@M16@@ parcourt les sous-@@M17@@-alg\`ebres de type fini de @@M18@@. En vertu de 3.1, le morphisme en question provient d'un morphisme @@M19@@. On peut donc supposer @@M20@@ de type fini sur @@M21@@.

Le morphisme @@M22@@ correspond \`a un morphisme f.p.p.f. @@M23@@ et un faisceau inversible @@M24@@ sur @@M25@@. Par descente f.p.p.f., on peut remplacer @@M26@@ par @@M27@@. Apr\`es le changement de base @@M28@@, compte tenu de la compatibilit\'e de la formation de @@M29@@ avec tout changement de base (4.1) on est donc ramen\'e dans (i) et (ii) \`a \'etudier le cas d'une section @@M30@@, d\'efinie par un faisceau inversible @@M31@@ sur @@M32@@.

On doit prouver que l'ensemble @@M33@@ des @@M34@@ tels que @@M35@@ soit @@M36@@-\'equivalent \`a @@M37@@ est ouvert, et aussi ferm\'e sous les conditions de (ii)\,; comme @@M38@@ est noeth\'erien, ceci veut dire trois choses~: pour un @@M39@@,
\begin{enumerate}[label=(\alph*)]
\item toute g\'en\'erisation de @@M40@@ reste dans @@M41@@.
\item toute sp\'ecialisation de @@M42@@, voisine de @@M43@@, reste dans @@M44@@.
\item sous les conditions de (ii), toute sp\'ecialisation de @@M45@@ reste dans @@M46@@.
\end{enumerate}

